\section*{Response to the reviewers}

We thank both reviewers for their careful reading of the manuscript and for their constructive comments. We have revised the manuscript accordingly, with the main changes including: (i) expanded discussion of film drops and recent fluid-mechanical perspectives, (ii) clarification of remote/open-ocean terminology, (iii) additional discussion of SSA modelling and parameterisation, (iv) revision of figure numbering and captions, and (v) clarification of number fluxes and particle size classes.

\section*{Response to Reviewer 1}

\subsection*{Comment 1}

\emph{Abstract: ``such as perfluoroalkyl acids'' may be a little too specific for the abstract (no other chemical is mentioned here).}

\textbf{Response:}

I agree that specifically highlighting perfluoroalkyl acids in the abstract places too much emphasis on a single chemical class relative to the broader scope of the article. I have therefore revised the sentence to refer more generally to the transfer of surface-active pollutants through sea spray aerosol.

\textbf{Changes made:}
The abstract has been revised to replace ``including the transfer of surface-active pollutants such as perfluoroalkyl acids'' with more general wording referring to the transfer of surface-active pollutants.

\subsection*{Comment 2}

\emph{Figure 1: by curiosity: these photographs show on all of them the same three droplets at the exact same place which is weird as they should have evolved from one picture to the other. Does the author explain what they are? If they are not real it may not be the best illustration of early science results..?}

\textbf{Response:}

The original paper explicitly explains that the stationary circular features visible in the image sequence are droplets adhering to the inner surface of the glass observation window. These droplets were produced by jets from earlier bubble-bursting events and therefore remain fixed relative to the camera frame while the bubble cavity and jet shown in the sequence evolve. I completely agree with the reviewer that this detail may not be immediately obvious to modern readers and have therefore clarified it in the figure caption.

\textbf{Changes made:}

The caption of Figure 1 has been revised to note that the stationary droplets visible in the image sequence are droplets adhering to the inner surface of the glass observation window, as described in Kientzler et al. (1954).

\subsection*{Comment 3}

\emph{Line 239: `In remote regions'': how is remote regions defined? As this statement contradicts the O'Dowd et al. 2004 results, do you suggest that remote is in contrast to coastal? would `open ocean'' be more appropriate?}

\textbf{Response:}

I agree that the term `remote regions'' is potentially ambiguous and that `open ocean'' more accurately reflects the environments sampled during the NAAMES and ATom campaigns discussed in this section. I have therefore revised the text accordingly.

\textbf{Changes made:}

The phrase `in remote regions'' has been replaced with `in the open ocean'' in the paragraph discussing the NAAMES and ATom observations.

\subsection*{Comment 4}

\emph{Line 255--256: ``distinguishing between film and jet drops and showing that film drops, which are responsible for submicron aerosols'': the distinction between jet and film drops, leading to different aerosol size classes, was also an early finding from the last century on which many modern research rely. It may be worth mentioning earlier in the paper.}

\textbf{Response:}

I agree that the distinction between jet drops and film drops is a foundational concept in SSA research and that introducing it only in the later discussion of laboratory studies may understate its historical importance. I have therefore expanded the early historical section to briefly discuss film-drop production, its role in generating submicron SSA, and the enduring importance of the jet-drop/film-drop distinction for understanding aerosol size distributions. I have also taken this opportunity to cite a recent synthesis of the underlying fluid mechanics.

\textbf{Changes made:}

A new paragraph has been added to the end of the section ``Early Foundations: The Jet Drop Mechanism'' introducing film-drop formation, its role in producing submicron SSA, and its historical significance, with reference to recent fluid-mechanical reviews.

\subsection*{Comment 5}

\emph{Line 329: ``such as viscosity and surface tension,..'' also density and evaporation rates may be important parameters.}

\textbf{Response:}

I agree that seawater density is another important temperature-dependent property that can influence bubble dynamics and aerosol production. I have revised the text to include density alongside viscosity and surface tension. While evaporation influences the evolution of aerosol particles after emission, the discussion here focuses specifically on seawater properties that influence bubble dynamics and droplet production.

\textbf{Changes made:}

The discussion of temperature-dependent seawater properties has been revised to refer to ``viscosity, density, and surface tension'' rather than only viscosity and surface tension.

\subsection*{Comment 6}

\emph{Line 333: `particle size'': would `particle size class (sub- vs supermicron)'' be more appropriate?}

\textbf{Response:}

I agree that the distinction between submicron and supermicron particles more accurately reflects how temperature effects on SSA production are commonly discussed in the literature. I have therefore revised the wording accordingly.

\textbf{Changes made:}

The phrase ``varies strongly with particle size'' has been revised to emphasise differences between particle size classes, particularly submicron and supermicron SSA.

\subsection*{Comment 7}

\emph{Line 366 `aerosol composition'': aerosol composition and number, and line 381 `chemical partitioning'' and number fluxes, as number fluxes are climatically more important than aerosol chemical composition.}

\textbf{Response:}

I agree that particle number concentrations and number fluxes are central to many climatic effects of SSA, particularly aerosol--cloud interactions and cloud condensation nuclei production. I have therefore revised the concluding section to explicitly acknowledge both aerosol composition and particle number, as well as both chemical partitioning and number fluxes.

\textbf{Changes made:}

The concluding discussion has been revised to refer to `the size-resolved emission fluxes, particle number concentrations, and chemical composition of SSA'' and to `chemical partitioning, number fluxes, biological modulation, and ultimately climate-relevant aerosol properties''.

\subsection*{Comment 8}

\emph{Figure numbering are missing?}

\textbf{Response:}

I thank the reviewer for identifying this issue. The omission of figure numbers was an unintended formatting problem in the submitted manuscript. This has now been corrected, and figure numbering is displayed consistently throughout the revised manuscript.

\textbf{Changes made:}

Figure numbering has been corrected in the revised manuscript.


\section*{Response to Reviewer 2}

\subsection*{Comment 1}

\emph{Jet drops are not the only mechanisms of sea drop production; another one is film drops. I suggest the recent Annual Review of Fluid Mechanics by Luc Deike in 2022, in particular, Section 5, dealing with fluid mechanics aspects of sea-spray aerosol.}

\textbf{Response:}

I agree that film drops are a fundamental pathway for sea-spray aerosol production and that the distinction between jet drops and film drops should be introduced earlier in the manuscript. I have therefore revised the early historical section to explicitly discuss film-drop formation during rupture of the bubble cap, its importance for submicron SSA production, and the continued relevance of the jet-drop/film-drop distinction for understanding aerosol size distributions. I have also added the suggested reference to the recent review by Deike (2022), which provides a valuable synthesis of the fluid-mechanical aspects of sea-spray aerosol production.

\textbf{Changes made:}

A new paragraph has been added to the section ``Early Foundations: The Jet Drop Mechanism'' discussing film-drop production, the different particle size classes associated with jet and film drops, and the relevance of recent fluid-mechanical work, including Deike (2022).

\subsection*{Comment 2}

\emph{The work is more descriptive than critical. It would be beneficial to identify where the field actually agrees, where it does not, if there are different theories of some aspects, and which uncertainties are most crucial for models and climate applications.}

\textbf{Response:}

I would like to thank the reviewer for this important comment. I agree that a perspective article should not only describe the historical development of the field but also identify areas of consensus, disagreement, and continuing uncertainty. I have therefore revised the manuscript to make the critical synthesis more explicit. In particular, I now more clearly distinguish between areas where the field broadly agrees, such as the importance of bubble bursting for SSA production and the occurrence of size-dependent organic enrichment, and areas where substantial uncertainty remains, including the extent to which short-term biological variability controls SSA composition across different oceanic regimes. I have also strengthened the concluding discussion to emphasise the uncertainties most relevant for models and climate applications, particularly size-resolved number fluxes, chemical composition, particle activation behaviour, and their dependence on physical, biological, and environmental conditions.

\textbf{Changes made:}

The discussion of organic enrichment and biological control has been revised to more explicitly identify agreement and disagreement within the field. The concluding section has also been revised to more clearly identify the uncertainties most important for models and climate applications, including size-resolved number fluxes, composition, and activation behaviour.

\subsection*{Comment 3}

\emph{An explicit equation/formulation of the parametrization described in Section 3 can be provided.}

\textbf{Response:}

I agree that including the original formulation helps illustrate why the work of Monahan and colleagues represented such an important advance for the field. I have therefore added the historical parameterisation introduced by Monahan et al. (1986), including the empirical whitecap-coverage relationship that linked wind speed to aerosol production. This addition more clearly demonstrates how physical observations of whitecaps were translated into quantitative source functions suitable for regional and global modelling studies.

\textbf{Changes made:}

Section 3 (``The Whitecap Era: Empirical Parameterisations'') now includes the original Monahan et al. (1986) formulation for SSA production and the associated whitecap-coverage parameterisation, together with a brief discussion of their significance for climate-model applications.

\subsection*{Comment 4}

\emph{The Figures are mentioned in the text as Figure 1 and Figure 2, but the Figure labels do not contain a number.}

\textbf{Response:}

I agree that figure numbering should be displayed consistently between the text and figure captions. The figure formatting has been revised to ensure that figure numbers are clearly shown in the captions and correspond to the numbering used in the main text.

\textbf{Changes made:}

The formatting of Figures 1 and 2 has been corrected so that figure numbers are displayed consistently in both the figure captions and the manuscript text.

\subsection*{Comment 5}

\emph{Figure 2 is directly extracted from the original paper, but the quality is not great. I suggest redoing the Figure by digitizing the data.}

\textbf{Response:}

I agree that the quality of the reproduced historical figure was limited by the resolution and reproduction quality of the original publication. I have therefore digitised the data from the original figure and produced a redrawn version that preserves the original information content while substantially improving clarity and readability. Appropriate attribution to the original \textit{Tellus} publication has been retained in the revised figure caption.

\textbf{Changes made:}

Figure 2 has been replaced with a redrawn version based on data digitised from Figure 4 of \citet{Blanchard1957}. The figure caption has been updated accordingly.
